\section{Lessons for SE Teams and Threats to Validity}
\label{sec:threats}
\subsection{Practical Lessons}
\noindent\textbf{Context dilution.} C2 outperforms C3 with both investigators. A relevant explanation may connect an assumption to a current-code obligation more directly than combined structural and textual history. Context dilution is a plausible explanation, not an isolated causal finding, since context length and attention are uncontrolled. Teams should retrieve history to resolve a specific hypothesis rather than maximize context volume.

\noindent\textbf{Verifiable boundary obligations.} SAST supplies reproducible entry points, while patch-selected symbols expose changes without warnings. A missing sibling guard or recorded trust assumption becomes actionable when translated into a check of current inputs, consumers, or object lifecycle. Preserve the link from the historical observation through the hypothesis to its validating evidence.

\noindent\textbf{Validation and maintenance gates.} Reserve final validation capacity, retain rejected patches and reasons, and reuse results only for identical candidates. Recognized terminal summaries distinguish completed checks from interrupted tests. Persistent memory carries the accepted invariant and independently owned check into maintenance, where the gate returns concrete feedback or preserves the secure start. Report feature delivery and fallback preservation separately.

\subsection{Threats to Validity and Scope}
\emph{Construct and internal validity.} Reference-patch overlap measures review opportunity, with broad classes and incidental edits potentially overstating semantic localization. Task tests cover specified contracts, not whole-repository security or unrelated alerts. Test adequacy and overfitting remain limits~\cite{repairsurvey,overfitting}, and recorded rationale need not faithfully describe internal reasoning~\cite{turpin,lanham}. Adaptive development changes several mechanisms together. C0 receives patch-selected targets, C1--C3 add audit cues without a cue-only control, and temporal evidence excludes upstream ancestry. Resumed execution retains earlier revisions, with paired original/candidate evaluations after environment recovery. The previously protected stratum is separately reported but is not an untouched prospective replication.

\emph{External and domain-study validity.} Selection favors Python outputs with publisher-reported functional success. Missing complete outcomes for 42 admissions may be nonrandom, and shared tasks create dependence addressed by task-cluster intervals. Few conversions limit precision. The fifteen semantic reviews are purposive and single-reviewer. Five component repairs in six structural-biology repositories are researcher-run with industry input, not independently confirmed CVEs or a practitioner deployment. Memory+Gate includes a retry absent from No Memory, so its gain cannot be attributed to memory alone. Source/check hashes establish traceability rather than semantic correctness. Broader repository coverage and independent practitioner validation remain future work.

\section{Related Work}
SWE-bench measures issue resolution~\cite{swebench}, whereas SusVibes separates functional and security outcomes~\cite{susvibes,susvibesartifact}. CVEfixes links vulnerabilities to fixing histories, and Vul4J provides reproducible vulnerabilities and triggering tests~\cite{cvefixes,vul4j}. SWE-agent studies repository interaction interfaces~\cite{sweagent}, while ReAct supplies interleaved explanation/action traces~\cite{react}. For security investigation, PAGENT combines static guidance with dynamic profiling~\cite{pagent}, IRIS infers taint specifications and filters candidate paths~\cite{iris}, LLMDFA couples model reasoning with parsing and feasibility checks~\cite{llmdfa}, and RepoAudit explores data flows using memory and validation~\cite{repoaudit}. Our comparison asks what development evidence adds to a common current-source repair protocol.

Historical analysis operates on complementary units. ViolationTracker reconstructs static-warning histories~\cite{violationtracker}, CoLeFunDa represents function changes to identify silent fixes~\cite{colefunda}, KNighter synthesizes checkers from historical patches~\cite{knighter}, and FixMiner extracts reusable repair patterns~\cite{fixminer}. GumTree provides structural edit matching~\cite{gumtree,gumtree3}. Analyzer-guided repair includes InferFix~\cite{inferfix}, while ChatRepair uses conversational test feedback~\cite{chatrepair} and RepairAgent autonomously inspects, edits, and tests code~\cite{repairagent}. These mechanisms motivate tracking losses between available evidence, supported diagnosis, and retained repair.

For subsequent maintenance, Reflexion retains verbal feedback in episodic memory~\cite{reflexion}, Self-Refine iterates through model-generated feedback~\cite{selfrefine}, and POSTURA describes a persistent security graph~\cite{postura}. The exploratory \emph{Undoing the Fix} report studies deliberately elicited regressions under an AI-control setting~\cite{undoingfix}. Our benign feature-evolution study instead connects an earlier repair invariant to independently owned executable checks, distinguishing first proposals, feedback-assisted acceptance, and secure fallback.

\section{Conclusion}
\tracer{} links change-informed planning, temporal AST evolution, recorded rationale, and guarded repair to persistent maintenance evidence. Across 3,096 cells, C2 achieves the highest shared-cohort SecPass counts with both investigators, reaching 76/345 with DeepSeek and 75/345 with Gemini from 68 original passes. Stage-level losses explain why broad review coverage produces modest repair gains. In structural-biology components, memory with a one-retry gate completes five feature evolutions while preserving saved boundaries, compared with three under unguarded evolution. The portable skill and archived evidence support auditable security review and maintenance.

\section*{Data Availability}
The skill, analysis helpers, derived per-cell outcomes, and tables are archived on Zenodo~\cite{tracerartifact} and maintained at \url{https://github.com/pkrajput/TRACER_skill}. SusVibes provides the benchmark inputs~\cite{susvibesartifact}. The archive supports outcome recomputation, but excludes raw investigator transcripts, complete evaluator logs, and container exports.

\section*{Acknowledgements}
We thank Google for Gemini API credits and Nostrum Biodiscovery for verification assistance and feedback on the work's practical relevance.
